%=============================================================================
\section{Modern Control Theory}
%=============================================================================

Classical control (PID, frequency response, root locus) got us far. It can handle most single-input, single-output systems beautifully. But some problems demand more:

\begin{itemize}
    \item Your holonomic chassis has four motors that must be coordinated simultaneously (MIMO).
    \item You want \textit{optimal} performance, not just ``good enough'' (optimal control).
    \item You need to estimate quantities you cannot measure directly, like velocity from position data or tilt angle from a noisy IMU (state estimation).
    \item You want to explicitly handle constraints like maximum motor current or joint angle limits (constrained optimization).
\end{itemize}

This is where modern control theory shines. It works directly with the state-space representation and uses linear algebra as its primary tool.

\subsection{Relationship Between State-Space and Transfer Functions}

These are two representations of the same system, and you can convert between them.

\textbf{State-space to transfer function:}
\begin{equation}
G(s) = C(sI - A)^{-1}B + D
\end{equation}

Taking the Laplace transform of $\dot{\mathbf{x}} = A\mathbf{x} + B\mathbf{u}$ (with zero initial conditions) gives $s\mathbf{X}(s) = A\mathbf{X}(s) + B\mathbf{U}(s)$, hence $\mathbf{X}(s) = (sI - A)^{-1}B\mathbf{U}(s)$. Substituting into $\mathbf{Y}(s) = C\mathbf{X}(s) + D\mathbf{U}(s)$ yields $G(s) = C(sI-A)^{-1}B + D$.

\textbf{Transfer function to state-space:} Not unique---many different state-space realizations can give the same transfer function. Common forms include \textit{controllable canonical form} (the denominator coefficients appear directly in the last row of $A$, making it easy to verify controllability) and \textit{observable canonical form} (the denominator coefficients appear in the last column of $A$, making observability transparent). In practice, MATLAB's \texttt{tf2ss} gives you one realization; which canonical form you use matters mainly for hand analysis.

\textbf{The key difference:} Transfer functions describe the input-output behavior. State-space describes the internal dynamics. A system can have identical transfer functions but different internal structures---the state-space representation captures this extra information.

\textbf{Why it matters:} Some internal modes might be unstable but invisible in the transfer function (if they are cancelled by zeros). State-space reveals them. This is why modern control prefers state-space---it sees everything.

\subsection{System Performance Metrics}

\subsubsection{Controllability and Observability}

These are two fundamental questions that you should ask before designing any controller or observer.

\textbf{Controllability:} ``Can I drive every state to any desired value using the available inputs?''

A system $(A, B)$ is controllable if the \textbf{controllability matrix} has full rank:
\begin{equation}
\mathcal{C} = \begin{bmatrix} B & AB & A^2B & \cdots & A^{n-1}B \end{bmatrix}, \quad \text{rank}(\mathcal{C}) = n
\end{equation}

\vspace{0.5em}
\begin{mdframed}[linewidth=1pt, innertopmargin=10pt, innerbottommargin=10pt]
\textbf{Derivation: Why the Controllability Matrix Works}

\vspace{0.3em}
Where does $\mathcal{C} = [B \;\; AB \;\; \cdots \;\; A^{n-1}B]$ come from? It is not pulled out of thin air---it follows directly from the state transition equation.

\textbf{Step 1: Iterative state propagation.} For the discrete-time system $\mathbf{x}[k+1] = A\mathbf{x}[k] + B\mathbf{u}[k]$, start from $\mathbf{x}[0] = \mathbf{0}$ and propagate forward:
\begin{align}
\mathbf{x}[1] &= B\mathbf{u}[0] \\
\mathbf{x}[2] &= A\mathbf{x}[1] + B\mathbf{u}[1] = AB\mathbf{u}[0] + B\mathbf{u}[1] \\
\mathbf{x}[3] &= A^2B\mathbf{u}[0] + AB\mathbf{u}[1] + B\mathbf{u}[2] \\
&\;\vdots \notag \\
\mathbf{x}[n] &= A^{n-1}B\mathbf{u}[0] + A^{n-2}B\mathbf{u}[1] + \cdots + B\mathbf{u}[n-1]
\end{align}

\textbf{Step 2: Rewrite as a matrix equation.} Stack the inputs into a vector:
\begin{equation}
\mathbf{x}[n] = \begin{bmatrix} A^{n-1}B & A^{n-2}B & \cdots & AB & B \end{bmatrix}
\begin{bmatrix} \mathbf{u}[0] \\ \mathbf{u}[1] \\ \vdots \\ \mathbf{u}[n-1] \end{bmatrix}
= \mathcal{C} \cdot \mathbf{u}_{\text{seq}}
\end{equation}

(The columns of $\mathcal{C}$ are just reordered---rank is unchanged.)

\textbf{Step 3: The rank condition.} We want to reach \textit{any} target state $\mathbf{x}[n] \in \mathbb{R}^n$. This requires the equation $\mathcal{C} \cdot \mathbf{u}_{\text{seq}} = \mathbf{x}[n]$ to have a solution for every right-hand side. By linear algebra, this is possible if and only if $\text{rank}(\mathcal{C}) = n$.

\textbf{Why only $n$ columns?} The Cayley--Hamilton theorem guarantees that $A^n$ can be written as a linear combination of $I, A, A^2, \ldots, A^{n-1}$. Therefore $A^n B$ is already in $\text{span}\{B, AB, \ldots, A^{n-1}B\}$, and higher powers add nothing new. You never need more than $n$ columns.

\textbf{Continuous-time version:} For $\dot{\mathbf{x}} = A\mathbf{x} + B\mathbf{u}$, the reachable set from $\mathbf{x}(0) = \mathbf{0}$ at time $t$ is:
\[
\mathbf{x}(t) = \int_0^{t} e^{A(t-\tau)} B \mathbf{u}(\tau) \, d\tau
\]
Expanding the matrix exponential $e^{A(t-\tau)} = \sum_{k=0}^{\infty} \frac{(t-\tau)^k}{k!} A^k$ and again applying Cayley--Hamilton to collapse powers $\geq n$, the reachable set is $\text{span}\{B, AB, \ldots, A^{n-1}B\}$---the same controllability matrix.
\end{mdframed}

\vspace{0.5em}
\begin{mdframed}[linewidth=1pt, innertopmargin=10pt, innerbottommargin=10pt]
\textbf{Derivation: Why the Observability Matrix Works}

\vspace{0.3em}
The observability matrix arises from the dual question: given outputs $\mathbf{y}[0], \mathbf{y}[1], \ldots$, can we recover the initial state $\mathbf{x}[0]$?

\textbf{Step 1: Collect output equations.} Assume $\mathbf{u} = \mathbf{0}$ (we can always subtract the known input contribution). The free response gives:
\begin{align}
\mathbf{y}[0] &= C\mathbf{x}[0] \\
\mathbf{y}[1] &= C\mathbf{x}[1] = CA\mathbf{x}[0] \\
\mathbf{y}[2] &= CA^2\mathbf{x}[0] \\
&\;\vdots \notag \\
\mathbf{y}[n-1] &= CA^{n-1}\mathbf{x}[0]
\end{align}

\textbf{Step 2: Stack into a matrix equation.}
\begin{equation}
\begin{bmatrix} \mathbf{y}[0] \\ \mathbf{y}[1] \\ \vdots \\ \mathbf{y}[n-1] \end{bmatrix}
= \underbrace{\begin{bmatrix} C \\ CA \\ CA^2 \\ \vdots \\ CA^{n-1} \end{bmatrix}}_{\mathcal{O}} \mathbf{x}[0]
\end{equation}

\textbf{Step 3: The rank condition.} To uniquely recover $\mathbf{x}[0]$ from the output sequence, the equation $\mathcal{O} \cdot \mathbf{x}[0] = \mathbf{y}_{\text{seq}}$ must have a unique solution. This requires $\text{rank}(\mathcal{O}) = n$ (full column rank).

\textbf{Why only $n$ rows?} Same reason as controllability---Cayley--Hamilton makes $CA^n$ a linear combination of $C, CA, \ldots, CA^{n-1}$, so additional measurements add no new information.

\textbf{The duality is now obvious:} Controllability asks ``can $\mathcal{C}$ map the input space \textit{onto} the state space?'' Observability asks ``can $\mathcal{O}$ map the state space \textit{into} distinguishable outputs?'' Transpose one and you get the other---which is exactly the mathematical statement of duality: $(A, B)$ controllable $\Leftrightarrow$ $(A^T, B^T)$ observable (with $B^T$ playing the role of $C$).
\end{mdframed}

\textbf{Physical interpretation:} If a state is uncontrollable, no input can influence it. For example, if your robot's chassis has a broken motor, you have lost controllability in certain directions.

\textbf{Worked example:} Consider a simplified 2D chassis with states $\mathbf{x} = \begin{bmatrix} v_x \\ v_y \end{bmatrix}$ (velocities in $x$ and $y$) and a single forward-facing motor:
\[
A = \begin{bmatrix} -1 & 0 \\ 0 & -1 \end{bmatrix}, \quad B = \begin{bmatrix} 1 \\ 0 \end{bmatrix}
\]
The controllability matrix: $\mathcal{C} = \begin{bmatrix} B & AB \end{bmatrix} = \begin{bmatrix} 1 & -1 \\ 0 & 0 \end{bmatrix}$. Rank = 1, but $n = 2$. \textbf{Not controllable!} The motor can only drive $v_x$; no input can change $v_y$. Physically obvious: a car cannot move sideways. Add a second motor facing sideways ($B = \begin{bmatrix} 1 & 0 \\ 0 & 1 \end{bmatrix}$) and controllability is restored.

\textbf{Observability:} ``Can I determine every state from the available outputs (measurements)?''

A system $(A, C)$ is observable if the \textbf{observability matrix} has full rank:
\begin{equation}
\mathcal{O} = \begin{bmatrix} C \\ CA \\ CA^2 \\ \vdots \\ CA^{n-1} \end{bmatrix}, \quad \text{rank}(\mathcal{O}) = n
\end{equation}

\textbf{Physical interpretation:} If a state is unobservable, you cannot infer it from measurements. For example, if you only measure motor position but not current, the current state might be unobservable (depending on the model).

\textbf{The controllability-observability duality:} A system $(A, B, C)$ is controllable if and only if the dual system $(A^T, C^T, B^T)$ is observable. This elegant symmetry means that techniques for controller design (pole placement) directly translate to observer design.

\textbf{Bottom line:} Always check controllability before designing a controller, and observability before designing an observer. If these conditions are not satisfied, no amount of mathematical sophistication will save you---you need to change the hardware (add sensors, fix actuators).

\subsubsection{Closed-Loop System Stability}

For the closed-loop state-space system with state feedback $\mathbf{u} = -K\mathbf{x}$:
\begin{equation}
\dot{\mathbf{x}} = (A - BK)\mathbf{x}
\end{equation}

The system is stable if and only if all eigenvalues of $(A - BK)$ have negative real parts. The gain matrix $K$ is what we design to place these eigenvalues where we want.

\textbf{Lyapunov stability:} A more general tool for stability analysis. A system is stable if there exists a positive definite function $V(\mathbf{x})$ (an ``energy-like'' function) that always decreases along trajectories ($\dot{V} < 0$). Finding such a $V$ proves stability; failing to find one does not prove instability (it might just mean you did not look hard enough).

For linear systems, the Lyapunov equation is:
\begin{equation}
A^T P + PA = -Q
\end{equation}

Choose $V(\mathbf{x}) = \mathbf{x}^T P \mathbf{x}$ as a candidate Lyapunov function. Its time derivative along trajectories of $\dot{\mathbf{x}} = A\mathbf{x}$ is $\dot{V} = \dot{\mathbf{x}}^T P \mathbf{x} + \mathbf{x}^T P \dot{\mathbf{x}} = \mathbf{x}^T(A^T P + PA)\mathbf{x}$. For $\dot{V} < 0$ (stability), we need $A^T P + PA = -Q$ for some positive definite $Q$.

If $Q > 0$ and the solution $P > 0$, the system is stable. MATLAB's \texttt{lyap()} solves this directly.

\subsection{Controller Design via State-Space}

\subsubsection{Bellman Equations and Dynamic Programming: The Engine Behind LQR and MPC}

Before we dive into LQR and MPC, we need to understand the mathematical engine that powers both: \textbf{dynamic programming (DP)} and the \textbf{Bellman equation}. If you have done competitive programming (USACO, Codeforces, LeetCode), you already know DP --- the control theory version is the same idea wearing a different hat.

\textbf{The Bellman Principle of Optimality:} \textit{An optimal policy has the property that, regardless of the initial state and initial decision, the remaining decisions must constitute an optimal policy with regard to the state resulting from the first decision.}

In plain English: if you know the optimal strategy from every intermediate state onward, you can figure out the optimal first move by trying all options and picking the one that minimizes (first-step cost) + (optimal cost-to-go from the resulting state).

\vspace{0.5em}
\begin{mdframed}[linewidth=1pt, innertopmargin=10pt, innerbottommargin=10pt]
\textbf{Warm-Up: Grid-World Shortest Path (DP You Already Know)}

\vspace{0.3em}
Consider a classic DP problem: you are on an $n \times n$ grid. You start at $(0, 0)$ and want to reach $(n{-}1, n{-}1)$. Each cell $(i, j)$ has a cost $c(i, j)$. You can move right or down. Find the minimum-cost path.

\textbf{State:} Your position $(i, j)$.

\textbf{Decision (control):} Move right or move down.

\textbf{Cost-to-go:} $V(i, j) = $ minimum total cost to reach the goal from $(i, j)$.

\textbf{Bellman equation:}
\[
V(i, j) = c(i, j) + \min\{V(i+1, j), \; V(i, j+1)\}
\]

\textbf{Boundary:} $V(n{-}1, n{-}1) = c(n{-}1, n{-}1)$.

\textbf{Solution:} Fill the table \textit{backward} from the goal to the start. At each cell, the optimal decision is whichever neighbor has the smaller cost-to-go. Total complexity: $O(n^2)$.

You have been doing this since your first DP problem. Now watch what happens when we change the vocabulary:

\begin{center}
\begin{tabular}{lll}
\toprule
\textbf{DP (programming)} & \textbf{Optimal control} & \textbf{What it means} \\
\midrule
State $(i, j)$ & State $\mathbf{x}_k$ & Where you are now \\
Decision (right/down) & Control input $\mathbf{u}_k$ & What you do next \\
Transition $(i,j) \to (i+1,j)$ & $\mathbf{x}_{k+1} = f(\mathbf{x}_k, \mathbf{u}_k)$ & Physics/dynamics \\
Step cost $c(i,j)$ & Stage cost $\ell(\mathbf{x}_k, \mathbf{u}_k)$ & Penalty for this step \\
Cost-to-go $V(i,j)$ & Value function $V_k(\mathbf{x})$ & Optimal future cost \\
Bellman recursion & Bellman equation & The DP update rule \\
Fill table backward & Backward Riccati recursion & How we solve it \\
Optimal action at $(i,j)$ & Optimal policy $\mathbf{u}_k^* = \pi(\mathbf{x}_k)$ & What to do at each state \\
\bottomrule
\end{tabular}
\end{center}

It is \textit{literally the same algorithm}. The only difference is that in control, the ``grid'' is continuous (states are real numbers, not integers), and the ``decisions'' are real-valued control inputs instead of discrete left/right choices.

\begin{figure}[H]
\centering
\includegraphics[width=\textwidth]{gridworld_dp.pdf}
\caption{Grid-world DP example ($5 \times 5$, moves: right or down). (a)~Each cell has a step cost $c(i,j)$. (b)~The cost-to-go $V(i,j)$ is filled \textit{backward} from the goal --- this is the DP table (= Riccati recursion in control). (c)~The optimal path (total cost = 13) follows the Bellman policy: at each cell, go toward the neighbor with the smaller $V$. Replace ``cell'' with ``state'' and ``right/down'' with ``control input'' and you have LQR.}
\end{figure}
\end{mdframed}

\textbf{The Bellman equation for optimal control (discrete time):}

Given a system $\mathbf{x}_{k+1} = f(\mathbf{x}_k, \mathbf{u}_k)$ with stage cost $\ell(\mathbf{x}_k, \mathbf{u}_k)$ and terminal cost $\ell_f(\mathbf{x}_N)$, the optimal cost-to-go (value function) satisfies:

\begin{equation}
V_k(\mathbf{x}) = \min_{\mathbf{u}} \left[ \ell(\mathbf{x}, \mathbf{u}) + V_{k+1}\left(f(\mathbf{x}, \mathbf{u})\right) \right]
\end{equation}

with boundary condition $V_N(\mathbf{x}) = \ell_f(\mathbf{x})$.

The optimal control at step $k$ is:
\begin{equation}
\mathbf{u}_k^*(\mathbf{x}) = \arg\min_{\mathbf{u}} \left[ \ell(\mathbf{x}, \mathbf{u}) + V_{k+1}\left(f(\mathbf{x}, \mathbf{u})\right) \right]
\end{equation}

\textbf{This is solved backward in time}, from $k = N$ to $k = 0$, just like filling the DP table from the goal to the start.

\vspace{0.5em}
\begin{mdframed}[linewidth=1pt, innertopmargin=10pt, innerbottommargin=10pt]
\textbf{From Grid-World to LQR: The ``Continuous Grid'' Problem}

\vspace{0.3em}
Now imagine the grid-world becomes continuous. Instead of cells, your ``position'' is a real-valued state $\mathbf{x} \in \mathbb{R}^n$. Instead of right/down, your ``move'' is a real-valued control $\mathbf{u} \in \mathbb{R}^m$. The transition is linear: $\mathbf{x}_{k+1} = A\mathbf{x}_k + B\mathbf{u}_k$. The step cost is quadratic: $\ell(\mathbf{x}, \mathbf{u}) = \mathbf{x}^T Q \mathbf{x} + \mathbf{u}^T R \mathbf{u}$.

Plug into the Bellman equation:
\[
V_k(\mathbf{x}) = \min_{\mathbf{u}} \left[ \mathbf{x}^T Q \mathbf{x} + \mathbf{u}^T R \mathbf{u} + V_{k+1}(A\mathbf{x} + B\mathbf{u}) \right]
\]

Here comes the magic: because the dynamics are linear and the cost is quadratic, the value function is \textbf{also quadratic}: $V_k(\mathbf{x}) = \mathbf{x}^T P_k \mathbf{x}$ for some matrix $P_k$. Substituting this ansatz and taking the derivative with respect to $\mathbf{u}$:

\[
\frac{\partial}{\partial \mathbf{u}} \left[ \mathbf{u}^T R \mathbf{u} + (A\mathbf{x} + B\mathbf{u})^T P_{k+1} (A\mathbf{x} + B\mathbf{u}) \right] = 0
\]

Solving gives:

Expand the Bellman equation: $V_k(\mathbf{x}) = \min_{\mathbf{u}}[\mathbf{x}^T Q \mathbf{x} + \mathbf{u}^T R \mathbf{u} + (A\mathbf{x}+B\mathbf{u})^T P_{k+1}(A\mathbf{x}+B\mathbf{u})]$. Differentiate the bracketed expression with respect to $\mathbf{u}$ and set to zero: $2R\mathbf{u} + 2B^T P_{k+1}(A\mathbf{x} + B\mathbf{u}) = 0$, which gives $\mathbf{u}^* = -(R + B^T P_{k+1} B)^{-1} B^T P_{k+1} A \mathbf{x}$. Substituting back yields the Riccati recursion for $P_k$.

\begin{align}
\mathbf{u}_k^* &= -\underbrace{(R + B^T P_{k+1} B)^{-1} B^T P_{k+1} A}_{K_k} \, \mathbf{x}_k \\
P_k &= Q + A^T P_{k+1} A - A^T P_{k+1} B (R + B^T P_{k+1} B)^{-1} B^T P_{k+1} A
\end{align}

This backward recursion for $P_k$ is the \textbf{discrete Riccati equation} --- exactly the ``fill the DP table backward'' step. When the horizon $N \to \infty$, $P_k$ converges to a constant $P$, and we get the steady-state \textbf{LQR} gain $K = (R + B^T P B)^{-1} B^T P A$.

\textbf{In the grid-world analogy:}
\begin{itemize}[nosep]
    \item $P_k$ is the cost-to-go matrix --- it tells you the ``remaining cost'' from any state, just like $V(i,j)$ in the grid.
    \item $K_k$ is the optimal policy at step $k$ --- it tells you what to do at each state, just like ``go right if $V(i+1,j) < V(i,j+1)$.''
    \item The Riccati recursion is the DP table fill --- same $O(N)$ backward pass.
\end{itemize}
\end{mdframed}

\textbf{Connection to MPC:} MPC is the same Bellman/DP framework, but with two additions:
\begin{enumerate}
    \item \textbf{Constraints} on states and inputs (like adding walls to the grid-world --- some cells are forbidden).
    \item \textbf{Receding horizon}: instead of solving once from $k = 0$ to $N$, re-solve at every timestep using the latest measurement as the new ``start position.''
\end{enumerate}

The constraints make the Bellman equation harder to solve (the $\min$ is now a constrained optimization), which is why MPC needs a QP solver or the ADMM approach of TinyMPC. But the underlying principle is identical: backward recursion to find the optimal cost-to-go, forward application of the optimal policy.

\textbf{Connection to reinforcement learning:} If you are familiar with RL, the Bellman equation is exactly the same as the \textbf{Bellman optimality equation} in RL ($V^*(s) = \max_a [r(s,a) + \gamma V^*(s')]$). LQR is the ``closed-form RL'' for linear-quadratic problems. When the dynamics are unknown or nonlinear, RL methods (Q-learning, policy gradient) approximate the solution that LQR computes exactly.

\subsubsection{LQR (Linear Quadratic Regulator)}

With the Bellman/DP framework fresh in mind, LQR is simply the steady-state solution to the linear-quadratic Bellman equation. Instead of manually placing poles, you specify a cost function that balances performance against effort:

\begin{equation}
J = \int_0^{\infty} \left( \mathbf{x}^T Q \mathbf{x} + \mathbf{u}^T R \mathbf{u} \right) dt
\end{equation}

The matrix $Q$ penalizes state deviations---larger $Q$ means more aggressive tracking and faster response---while $R$ penalizes control effort---larger $R$ means gentler inputs and energy saving. The balance between $Q$ and $R$ is the fundamental design knob of LQR.

The optimal state feedback gain is:
\begin{equation}
K = R^{-1} B^T P
\end{equation}

where $P$ is the solution to the \textbf{continuous algebraic Riccati equation (CARE)}:
\begin{equation}
A^T P + PA - PBR^{-1}B^T P + Q = 0
\end{equation}

Do not solve this by hand. Use \texttt{lqr(A, B, Q, R)} in MATLAB or \texttt{scipy.linalg.solve\_continuous\_are()} in Python.

\textbf{What is the Riccati equation actually computing?} The matrix $P$ represents the \textbf{optimal cost-to-go}: $P_{ij}$ tells you how much ``future cost'' is associated with the interaction between states $i$ and $j$. Think of it as a map of ``how expensive is it to be in each state?'' States that are hard to control (require lots of effort to bring back to zero) get large entries in $P$. The optimal gain $K = R^{-1}B^T P$ then says: ``Apply control effort proportional to how costly each state is.'' Expensive states get corrected aggressively; cheap states get gentle treatment.

\textbf{Tuning LQR.}\;The art is in choosing $Q$ and $R$. Start with $Q = I$ (identity matrix) and $R = \rho I$ where $\rho$ is a single scalar. A small $\rho$ gives aggressive, fast control that uses lots of actuator effort; a large $\rho$ is conservative and gentle on actuators. To penalize specific states more, increase the corresponding diagonal elements of $Q$---for example, if you care more about position than velocity, make $Q_{11} > Q_{22}$.

\textbf{Why LQR is great for competitions.}\;It provides guaranteed stability (for controllable systems) and guaranteed robustness margins---at least 6 dB gain margin and 60\textdegree{} phase margin for SISO systems. It handles MIMO systems effortlessly, and the $Q/R$ tuning is more intuitive than PID's three gains for multi-axis systems.

\textbf{Example:} Balancing a two-wheeled robot. The state is $\mathbf{x} = [\theta, \dot{\theta}, x, \dot{x}]^T$ (tilt angle, tilt rate, position, velocity). PID would need multiple nested loops; LQR handles it with a single $K$ matrix.

\textbf{Limitation:} LQR requires full state feedback ($\mathbf{u} = -K\mathbf{x}$), which means you need to know all states. If you cannot measure all states, you need an observer (see Section 6.4).

\subsubsection{A Brief Primer on Optimization and Convexity}

MPC solves an optimization problem at every time step. Before diving in, let us establish the minimum optimization vocabulary you need.

\textbf{The general optimization problem:}
\[
\min_{\mathbf{x}} \; f(\mathbf{x}) \quad \text{subject to} \quad g_i(\mathbf{x}) \leq 0, \;\; h_j(\mathbf{x}) = 0
\]
where $f$ is the \textbf{objective function} (cost), $g_i$ are \textbf{inequality constraints}, and $h_j$ are \textbf{equality constraints}. The optimizer searches for the decision variable $\mathbf{x}$ that minimizes $f$ without violating any constraint.

\textbf{What makes a problem convex?} A problem is convex if:
\begin{enumerate}[nosep]
    \item The objective $f$ is a \textbf{convex function}: for any two points $\mathbf{a}, \mathbf{b}$ and $0 \leq \lambda \leq 1$,
    \[
    f\bigl(\lambda \mathbf{a} + (1-\lambda)\mathbf{b}\bigr) \leq \lambda f(\mathbf{a}) + (1-\lambda) f(\mathbf{b})
    \]
    Geometrically: the function lies below any chord. Quadratic functions $f(\mathbf{x}) = \mathbf{x}^T Q \mathbf{x}$ with $Q \succeq 0$ (positive semidefinite) are convex.
    \item The \textbf{feasible set} (defined by the constraints) is convex: any line segment between two feasible points stays feasible. Linear inequalities ($A\mathbf{x} \leq \mathbf{b}$) always define convex sets.
\end{enumerate}

\textbf{Why convexity matters:} A convex problem has a single global minimum---no local minima traps. Any solver that finds a solution has found \textit{the} solution. Non-convex problems can have many local minima, and solvers may get stuck in a bad one.

\textbf{Quadratic Programs (QP):} The most common optimization problem in control. The objective is quadratic, the constraints are linear:
\[
\min_{\mathbf{x}} \; \tfrac{1}{2}\mathbf{x}^T H \mathbf{x} + \mathbf{c}^T \mathbf{x} \quad \text{s.t.} \quad A\mathbf{x} \leq \mathbf{b}
\]
With $H \succeq 0$, this is convex and can be solved efficiently even on embedded systems (TinyMPC uses this structure). LQR without constraints is an \textit{unconstrained} QP whose solution is the Riccati equation. MPC with box constraints on inputs and states is a \textit{constrained} QP.

\textbf{Connection to control:} Every optimal control problem we encounter in this chapter---LQR, MPC, Kalman filter (dual of LQR)---is a convex QP. This is not a coincidence: the linear dynamics and quadratic cost were chosen \textit{precisely because} they make the problem convex and therefore tractable in real time.

\subsubsection{MPC (Model Predictive Control)}

MPC is the most powerful and flexible controller in the modern control toolbox. It solves an optimization problem at every time step:

\textit{``Given where I am now, what sequence of inputs over the next $N$ steps minimizes my cost function while respecting all constraints?''}

\textbf{At each time step:}
\begin{enumerate}
    \item Measure (or estimate) the current state $\mathbf{x}[k]$.
    \item Solve the optimization problem:
    \begin{equation}
    \min_{\mathbf{u}[k], \ldots, \mathbf{u}[k+N-1]} \sum_{i=0}^{N-1} \left( \mathbf{x}[k+i]^T Q \mathbf{x}[k+i] + \mathbf{u}[k+i]^T R \mathbf{u}[k+i] \right) + \mathbf{x}[k+N]^T P_f \mathbf{x}[k+N]
    \end{equation}
    subject to:
    \begin{align}
    \mathbf{x}[k+i+1] &= A_d \mathbf{x}[k+i] + B_d \mathbf{u}[k+i] \\
    \mathbf{u}_{\min} &\leq \mathbf{u}[k+i] \leq \mathbf{u}_{\max} \\
    \mathbf{x}_{\min} &\leq \mathbf{x}[k+i] \leq \mathbf{x}_{\max}
    \end{align}
    \item Apply only the first input $\mathbf{u}[k]$.
    \item Repeat at the next time step (receding horizon).
\end{enumerate}

\textbf{Understanding the cost function terms:}
\begin{itemize}
    \item $\mathbf{x}^T Q \mathbf{x}$: \textbf{Stage cost on states.} Penalizes deviation from the desired state at each step. Larger $Q_{ii}$ means ``I care more about keeping state $i$ close to zero (or the reference).''
    \item $\mathbf{u}^T R \mathbf{u}$: \textbf{Stage cost on inputs.} Penalizes control effort. Larger $R$ means ``go easy on the actuators.'' This prevents the optimizer from finding mathematically optimal but physically violent solutions.
    \item $\mathbf{x}[k+N]^T P_f \mathbf{x}[k+N]$: \textbf{Terminal cost.} Penalizes the state at the end of the prediction horizon. Without this, MPC can ``cheat'' by pushing all the error to the end of the horizon (like a student who does nothing until the last day). Setting $P_f$ to the solution of the LQR Riccati equation is a common choice that guarantees stability.
    \item $N$ (\textbf{prediction horizon}): How many steps ahead MPC looks. Longer horizon = better anticipation but heavier computation. Too short = myopic behavior (MPC cannot ``see'' upcoming turns or obstacles). A good starting point: $N = 10$--$30$ steps.
\end{itemize}

\textbf{The receding horizon concept:} MPC computes an \textit{entire} optimal trajectory but only executes the first step, then replans. This sounds wasteful, but it is what makes MPC robust: at each step, it uses the latest measurement to correct for model errors and disturbances. It is like a chess player who plans 10 moves ahead but re-evaluates after each opponent's move.

\textbf{Why MPC is powerful:}
\begin{itemize}
    \item \textbf{Handles constraints explicitly.} Motor current limits, joint angle limits, velocity bounds---they are all part of the optimization.
    \item \textbf{Handles MIMO systems naturally.} Multiple inputs and outputs are just more variables in the optimization.
    \item \textbf{Preview capability.} If you know the future reference trajectory (e.g., a planned path), MPC can look ahead and act proactively.
    \item \textbf{Intuitive tuning.} Like LQR, you tune $Q$ and $R$. The prediction horizon $N$ is the main additional parameter.
\end{itemize}

\textbf{The catch:} MPC requires solving an optimization problem at every time step. For linear systems with quadratic costs (linear MPC), this is a quadratic program (QP) that can be solved efficiently. For nonlinear systems (nonlinear MPC), it is much harder.

\textbf{Practical considerations for competition robots:}
For simple systems (chassis velocity control), PID or LQR is usually sufficient---MPC is overkill. But for trajectory tracking with constraints (autonomous navigation, robotic arms), MPC is worth the additional implementation effort. Libraries like OSQP, qpOASES, and acados can solve the QP fast enough for real-time embedded use. When even that is too expensive, \textit{explicit MPC} precomputes the solution offline as a lookup table indexed by the current state.

\begin{figure}[H]
\centering
\includegraphics[width=\textwidth]{mpc_horizon.pdf}
\caption{MPC receding horizon illustrated on a position-tracking problem. At each timestep, MPC plans an optimal trajectory over the next $N=10$ steps (green), but only applies the \textbf{first} control action (circled). Then it re-measures, replans, and repeats. The prediction window slides forward with each step---this ``plan ahead, execute one, replan'' cycle is what makes MPC robust to model errors and disturbances.}
\end{figure}

\subsubsection{TinyMPC: Riccati-Accelerated MPC for Embedded Systems}

Standard MPC solves a full QP (quadratic program) at every timestep. For a system with $n$ states, $m$ inputs, and horizon $N$, this means solving a problem with $N(n+m)$ decision variables --- feasible on a laptop but brutal on a microcontroller. \textbf{TinyMPC} (Nguyen et al., 2024) is a breakthrough that makes MPC run on resource-constrained embedded systems (Teensy, STM32, even Arduino) by exploiting the Riccati equation structure.

\textbf{The key insight:} For an unconstrained linear MPC, the solution is just a time-varying LQR --- you can compute it via a backward Riccati recursion. TinyMPC precomputes this Riccati recursion \textbf{offline}, then handles constraints \textbf{online} with a lightweight ADMM (Alternating Direction Method of Multipliers) loop.

\textbf{Algorithm structure:}

\textit{Step 1 --- Offline precomputation (once, or when model changes):}

Run the backward Riccati recursion from the terminal cost $P_N = Q_f$ to compute the optimal feedback gains at each step:
\begin{align}
S_k &= R + B_d^T P_{k+1} B_d \\
K_k &= S_k^{-1} B_d^T P_{k+1} A_d \\
P_k &= Q + A_d^T P_{k+1} A_d - A_d^T P_{k+1} B_d K_k
\end{align}

This gives you $N$ feedback gain matrices $K_0, K_1, \ldots, K_{N-1}$ and $N+1$ cost-to-go matrices $P_0, P_1, \ldots, P_N$. These are stored in memory and \textit{never recomputed} (unless the model or cost matrices change).

\textit{Step 2 --- Online solve (every control cycle):}

Given the current state $\mathbf{x}_0$:
\begin{enumerate}
    \item \textbf{Forward rollout:} Compute the unconstrained optimal trajectory using the precomputed gains: $\mathbf{u}_k = -K_k (\mathbf{x}_k - \mathbf{x}_k^{\text{ref}})$, then $\mathbf{x}_{k+1} = A_d \mathbf{x}_k + B_d \mathbf{u}_k$.
    \item \textbf{ADMM projection:} If any $\mathbf{u}_k$ or $\mathbf{x}_k$ violates constraints, run a few ADMM iterations (typically 5--20) to find the nearest feasible solution. Each ADMM iteration is just:
    \begin{itemize}[nosep]
        \item A backward Riccati-like pass (using the \textit{precomputed} $K_k$, so it is cheap).
        \item A projection step: $z_k = \text{clamp}(u_k + \lambda_k, u_{\min}, u_{\max})$ --- element-wise clamping.
        \item A dual update: $\lambda_k \leftarrow \lambda_k + u_k - z_k$.
    \end{itemize}
    \item \textbf{Apply first input:} Send $\mathbf{u}_0$ to the actuator. Discard the rest (receding horizon).
\end{enumerate}

\textbf{Why this is so fast:}
\begin{itemize}
    \item The Riccati recursion (the expensive part) is done \textbf{offline}. Online work is just matrix multiplies and clamping.
    \item No QP solver library needed. No matrix factorization online. Just multiply-add and clamp.
    \item Each ADMM iteration is $O(N \cdot n^2)$ --- linear in horizon, quadratic in state dimension. For a 4-state system with $N = 10$: about 1600 multiply-adds per iteration. At 10 iterations, that is 16,000 multiply-adds --- an STM32F4 at 168 MHz finishes this in $< 100\,\mu$s.
    \item Memory usage: $N \times (n \times n + m \times n)$ floats for the precomputed gains. For 4 states, 2 inputs, $N = 10$: about 240 floats = 960 bytes.
\end{itemize}

\textbf{Comparison with standard QP-based MPC:}

\begin{center}
\begin{tabular}{lcc}
\toprule
 & \textbf{Standard MPC (OSQP)} & \textbf{TinyMPC} \\
\midrule
Online computation & Full QP solve & Riccati rollout + ADMM \\
Typical solve time (STM32) & 1--10 ms & 50--200 $\mu$s \\
Memory & KBs (sparse matrices) & $<$ 1 KB for small systems \\
Dependencies & QP solver library & None (just matrix ops) \\
Constraint handling & Exact (to tolerance) & Approximate (ADMM) \\
Warm starting & Possible & Natural (previous solution) \\
\bottomrule
\end{tabular}
\end{center}

\textbf{When to use TinyMPC:}
\begin{itemize}
    \item You need MPC on a microcontroller with limited RAM/CPU.
    \item Your model is linear (or can be linearized at each step for nonlinear systems).
    \item Hard constraints on inputs are important (motor current limits, joint angle limits).
    \item Horizon is short to moderate ($N \leq 30$).
\end{itemize}

\textbf{When NOT to use TinyMPC:}
\begin{itemize}
    \item Complex nonlinear models (use acados or CasADi on a more powerful computer).
    \item Very long horizons ($N > 50$) with many state constraints.
    \item If your problem is unconstrained --- just use LQR, it is even simpler.
\end{itemize}

\textbf{Reference:} The full TinyMPC paper and open-source implementation are at \texttt{tinympc.org}. The C implementation fits in under 1000 lines and runs on ARM Cortex-M microcontrollers. Our C++ module in the Appendix follows the same Riccati + ADMM structure.

\subsection{Observer Design via State-Space}

LQR (and state feedback in general) requires knowing all states. But in practice, you can only measure some of them. An \textbf{observer} (or \textbf{estimator}) reconstructs the full state vector from the available measurements.

\subsubsection{Luenberger Observer}

The Luenberger observer is a copy of the system model, with a correction term that uses the measurement error to drive the estimate toward the true state.

\begin{equation}
\dot{\hat{\mathbf{x}}} = A\hat{\mathbf{x}} + B\mathbf{u} + L(\mathbf{y} - C\hat{\mathbf{x}})
\end{equation}

where:
\begin{itemize}[nosep]
    \item $\hat{\mathbf{x}}$ is the estimated state.
    \item $\mathbf{y} = C\mathbf{x}$ is the actual measurement.
    \item $L$ is the observer gain matrix (we design this).
\end{itemize}

The estimation error $\mathbf{e} = \mathbf{x} - \hat{\mathbf{x}}$ evolves as:
\begin{equation}
\dot{\mathbf{e}} = (A - LC)\mathbf{e}
\end{equation}

Subtract the observer equation $\dot{\hat{\mathbf{x}}} = A\hat{\mathbf{x}} + B\mathbf{u} + L(\mathbf{y} - C\hat{\mathbf{x}})$ from the true system $\dot{\mathbf{x}} = A\mathbf{x} + B\mathbf{u}$. The $B\mathbf{u}$ terms cancel, and substituting $\mathbf{y} = C\mathbf{x}$ gives $\dot{\mathbf{e}} = A\mathbf{e} - LC\mathbf{e} = (A - LC)\mathbf{e}$.

If we choose $L$ such that all eigenvalues of $(A - LC)$ have negative real parts, the estimation error decays to zero. The observer ``catches up'' to the true state.

\textbf{Design rule of thumb:} Place the observer poles 2--5$\times$ faster than the controller poles. The observer should converge faster than the controller acts, so the controller always has a good state estimate.

\textbf{Separation principle:} You can design the controller gain $K$ and observer gain $L$ independently. The combined controller-observer system is stable if each part is stable individually. This is a beautiful theoretical result that simplifies design enormously.

\textbf{When to use Luenberger:} When your model is accurate and the noise characteristics are not well known. It is simple, deterministic, and easy to debug.

\subsubsection{The Deep Duality: Controller and Observer as Mirror Images}

This is one of the most elegant results in all of control theory, and understanding it deeply will make you a better engineer. Let us lay out the parallel structure explicitly.

\textbf{The controller problem:} Given the system $\dot{\mathbf{x}} = A\mathbf{x} + B\mathbf{u}$, find a gain $K$ such that $\mathbf{u} = -K\mathbf{x}$ makes the closed-loop system $\dot{\mathbf{x}} = (A - BK)\mathbf{x}$ stable and well-behaved.

\textbf{The observer problem:} Given the system $\dot{\mathbf{x}} = A\mathbf{x} + B\mathbf{u}$, $\mathbf{y} = C\mathbf{x}$, find a gain $L$ such that the estimation error $\dot{\mathbf{e}} = (A - LC)\mathbf{e}$ decays to zero.

Now stare at the two closed-loop matrices side by side:

\begin{center}
\begin{tabular}{lcc}
\toprule
 & \textbf{Controller} & \textbf{Observer} \\
\midrule
Closed-loop matrix & $A - BK$ & $A - LC$ \\
Gain to design & $K$ & $L$ \\
Required condition & Controllability $(A, B)$ & Observability $(A, C)$ \\
Design goal & Place eigenvalues of $A - BK$ & Place eigenvalues of $A - LC$ \\
Information flow & Input $\to$ states & Measurements $\to$ state estimates \\
Riccati equation & $A^TP + PA - PBR^{-1}B^TP + Q = 0$ & $AP_o + P_oA^T - P_oC^TR_n^{-1}CP_o + Q_n = 0$ \\
\bottomrule
\end{tabular}
\end{center}

\textbf{The mathematical duality:} If you transpose the observer Riccati equation, replacing $A \to A^T$, $C^T \to B$, $Q_n \to Q$, $R_n \to R$, you get exactly the controller Riccati equation. They are literally the same equation with transposed matrices. This means:

\begin{itemize}
    \item Any algorithm that designs a controller can design an observer (just transpose the inputs).
    \item If you can place poles for a controller using \texttt{place(A, B, poles)}, you can place observer poles using \texttt{place(A', C', poles)'} --- note the transposes.
    \item The LQR controller ($K$ from the control Riccati) has a dual: the \textbf{Kalman filter} ($L$ from the observer Riccati). LQR is the optimal controller; Kalman is the optimal observer. They are dual problems.
\end{itemize}

\textbf{A physical analogy:} Think of the controller as a \textit{speaker} and the observer as a \textit{microphone}. The speaker (controller) pushes energy into the system through the inputs $B$ to shape the state trajectory. The microphone (observer) listens to the system through the outputs $C$ to reconstruct what is happening inside. Controllability asks: ``Can the speaker reach every corner of the room?'' Observability asks: ``Can the microphone hear every corner of the room?'' They are symmetric questions about the same room (system).

\textbf{Why this matters in practice:} If you have already tuned a state feedback controller and it works well, you have a strong intuition for how pole placement affects performance. That \textit{exact same intuition} applies when designing an observer. Faster poles = faster convergence but more noise sensitivity. Slower poles = smoother estimates but slower tracking. You are turning the same knob, just on the other side of the mirror.

\subsubsection{Kalman Filter}

The Kalman filter is the optimal observer for linear systems with Gaussian noise. It is the Luenberger observer's probabilistic cousin---instead of placing observer poles manually, it computes the optimal gain $L$ based on the noise statistics.

\textbf{System model with noise:}
\begin{align}
\mathbf{x}[k+1] &= A_d \mathbf{x}[k] + B_d \mathbf{u}[k] + \mathbf{w}[k], \quad \mathbf{w} \sim \mathcal{N}(0, Q_n) \\
\mathbf{y}[k] &= C \mathbf{x}[k] + \mathbf{v}[k], \quad \mathbf{v} \sim \mathcal{N}(0, R_n)
\end{align}

where $\mathbf{w}$ is process noise (model uncertainty) and $\mathbf{v}$ is measurement noise. $Q_n$ and $R_n$ are their covariance matrices.

\textbf{The Kalman filter algorithm (discrete):}

\textit{Predict:}
\begin{align}
\hat{\mathbf{x}}[k|k-1] &= A_d \hat{\mathbf{x}}[k-1|k-1] + B_d \mathbf{u}[k-1] \\
P[k|k-1] &= A_d P[k-1|k-1] A_d^T + Q_n
\end{align}

\textit{Update:}
\begin{align}
K_k &= P[k|k-1] C^T (C P[k|k-1] C^T + R_n)^{-1} \\
\hat{\mathbf{x}}[k|k] &= \hat{\mathbf{x}}[k|k-1] + K_k (\mathbf{y}[k] - C \hat{\mathbf{x}}[k|k-1]) \\
P[k|k] &= (I - K_k C) P[k|k-1]
\end{align}

\textbf{Intuition:} The Kalman gain $K_k$ automatically balances trust between the model prediction and the measurement:
\begin{itemize}
    \item If process noise is large ($Q_n$ big) relative to measurement noise ($R_n$ small), the filter trusts the measurement more.
    \item If measurement noise is large relative to process noise, the filter trusts the model more.
\end{itemize}

\textbf{Tuning the Kalman filter} comes down to choosing $Q_n$ and $R_n$:
\begin{itemize}
    \item $R_n$ can often be measured directly (record sensor data when the system is stationary, compute the variance).
    \item $Q_n$ is harder---it represents model uncertainty. Start with small values and increase if the filter is too sluggish (not tracking fast enough).
    \item A common cheat: treat $Q_n$ and $R_n$ as tuning knobs. Increase $Q_n / R_n$ ratio for faster tracking (more noise). Decrease for smoother estimates (more lag). Sounds familiar? It is the same trade-off as low-pass filter cutoff frequency, just expressed probabilistically.
\end{itemize}

\textbf{Why the Kalman filter is everywhere in robotics:}
\begin{itemize}
    \item \textbf{Sensor fusion:} Combine IMU, encoder, and vision data into one coherent state estimate.
    \item \textbf{Velocity estimation:} Get smooth velocity from noisy position measurements without a pure differentiator.
    \item \textbf{Optimal in a well-defined sense:} It minimizes the mean-squared estimation error for linear Gaussian systems.
\end{itemize}

\textbf{A practical note:} The Kalman filter needs a matrix inverse at each step (the $K_k$ computation). For small state dimensions ($n \leq 6$), this is trivial. For larger systems, use efficient implementations (e.g., Cholesky decomposition). On an STM32F4, a 6-state EKF runs comfortably at 1 kHz.

\subsubsection{Extended Kalman Filter (EKF)}

The standard Kalman filter assumes linear dynamics ($\mathbf{x}_{k+1} = A\mathbf{x}_k + B\mathbf{u}_k$) and linear measurements ($\mathbf{y}_k = C\mathbf{x}_k$). But real robotics systems are \textbf{nonlinear}: the relationship between Euler angles and angular rates involves trigonometric functions, the accelerometer measurement model depends on the sine/cosine of the tilt angles, and so on.

The \textbf{Extended Kalman Filter} handles nonlinearity by \textit{linearizing} the system at each timestep around the current estimate. The idea is simple: at each step, pretend the system is locally linear, apply the standard Kalman filter, then re-linearize at the new estimate.

\textbf{Nonlinear system model:}
\begin{align}
\mathbf{x}_{k+1} &= f(\mathbf{x}_k, \mathbf{u}_k) + \mathbf{w}_k \quad \text{(nonlinear dynamics)} \\
\mathbf{y}_k &= h(\mathbf{x}_k) + \mathbf{v}_k \quad \text{(nonlinear measurement)}
\end{align}

\textbf{The Jacobian matrix --- what it is and why we need it.}\;The Kalman filter's covariance propagation ($P_{k|k-1} = A P_{k-1|k-1} A^T + Q$) requires a \emph{matrix} $A$. But for a nonlinear system, there is no single matrix $A$---the ``slope'' of $f$ changes everywhere. The solution is to use the \textbf{Jacobian matrix}, which is simply the multivariable generalization of the derivative.

Recall from calculus: for a scalar function $y = f(x)$, the derivative $f'(x_0)$ tells you the best linear approximation near $x_0$: $f(x) \approx f(x_0) + f'(x_0)(x - x_0)$. The Jacobian does the same thing for a \emph{vector-valued} function of \emph{multiple variables}. If $\mathbf{f}: \mathbb{R}^n \to \mathbb{R}^m$, the Jacobian is the $m \times n$ matrix of all partial derivatives:
\[
F = \frac{\partial \mathbf{f}}{\partial \mathbf{x}} = \begin{bmatrix}
\frac{\partial f_1}{\partial x_1} & \cdots & \frac{\partial f_1}{\partial x_n} \\
\vdots & \ddots & \vdots \\
\frac{\partial f_m}{\partial x_1} & \cdots & \frac{\partial f_m}{\partial x_n}
\end{bmatrix}.
\]
Each row tells you how one output component changes with respect to all inputs. Evaluated at a specific point $\mathbf{x}_0$, the Jacobian gives the best local linear approximation: $\mathbf{f}(\mathbf{x}) \approx \mathbf{f}(\mathbf{x}_0) + F(\mathbf{x} - \mathbf{x}_0)$---exactly the ``pretend the system is locally linear'' step the EKF needs.

\textbf{EKF algorithm:}

\textit{Predict:}
\begin{align}
\hat{\mathbf{x}}_{k|k-1} &= f(\hat{\mathbf{x}}_{k-1|k-1}, \mathbf{u}_{k-1}) \\
F_k &= \left.\frac{\partial f}{\partial \mathbf{x}}\right|_{\hat{\mathbf{x}}_{k-1|k-1}} \quad \text{(Jacobian of dynamics)} \\
P_{k|k-1} &= F_k P_{k-1|k-1} F_k^T + Q
\end{align}

\textit{Update:}
\begin{align}
H_k &= \left.\frac{\partial h}{\partial \mathbf{x}}\right|_{\hat{\mathbf{x}}_{k|k-1}} \quad \text{(Jacobian of measurement)} \\
K_k &= P_{k|k-1} H_k^T (H_k P_{k|k-1} H_k^T + R)^{-1} \\
\hat{\mathbf{x}}_{k|k} &= \hat{\mathbf{x}}_{k|k-1} + K_k \left(\mathbf{y}_k - h(\hat{\mathbf{x}}_{k|k-1})\right) \\
P_{k|k} &= (I - K_k H_k) P_{k|k-1}
\end{align}

\textbf{The only difference from the standard Kalman filter:}
\begin{enumerate}
    \item The prediction uses the \textit{nonlinear} function $f(\cdot)$ instead of the matrix $A$.
    \item The measurement innovation uses the \textit{nonlinear} function $h(\cdot)$ instead of the matrix $C$.
    \item The covariance propagation uses the \textit{Jacobians} $F_k$ and $H_k$ (evaluated at the current estimate) instead of the constant matrices $A$ and $C$.
\end{enumerate}

Everything else --- the Kalman gain formula, the covariance update, the predict-update cycle --- is \textit{identical}.

\vspace{0.5em}
\begin{mdframed}[linewidth=1pt, innertopmargin=10pt, innerbottommargin=10pt]
\textbf{Case Study: 6-DOF IMU Pose Estimation with EKF}

\vspace{0.3em}
This is the single most common EKF application in competition robotics: estimating the 3D orientation (roll, pitch, yaw) of your robot from a 6-axis IMU (3-axis accelerometer + 3-axis gyroscope).

\textbf{State vector:} $\mathbf{x} = [\phi, \theta, \psi, b_x, b_y, b_z]^T$
\begin{itemize}[nosep]
    \item $\phi$: roll, $\theta$: pitch, $\psi$: yaw (Euler angles in radians)
    \item $b_x, b_y, b_z$: gyroscope biases (rad/s) --- the EKF estimates these automatically
\end{itemize}

\textbf{Dynamics model $f(\mathbf{x}, \mathbf{u})$:} The angular rates from the gyroscope (after bias subtraction) drive the Euler angles via the kinematic equations:
\begin{align}
\omega_x &= \text{gyro}_x - b_x, \quad \omega_y = \text{gyro}_y - b_y, \quad \omega_z = \text{gyro}_z - b_z \\[3pt]
\dot{\phi} &= \omega_x + \sin\phi \tan\theta \cdot \omega_y + \cos\phi \tan\theta \cdot \omega_z \\
\dot{\theta} &= \cos\phi \cdot \omega_y - \sin\phi \cdot \omega_z \\
\dot{\psi} &= \frac{\sin\phi}{\cos\theta} \cdot \omega_y + \frac{\cos\phi}{\cos\theta} \cdot \omega_z \\[3pt]
\dot{b}_x &= 0, \quad \dot{b}_y = 0, \quad \dot{b}_z = 0 \quad \text{(biases modeled as random walk)}
\end{align}

The Jacobian $F = \frac{\partial f}{\partial \mathbf{x}}$ is a $6 \times 6$ matrix with trigonometric entries. It must be recomputed at every timestep (this is the ``extended'' part).

\textbf{Measurement model $h(\mathbf{x})$:} The accelerometer measures gravity direction, which gives us roll and pitch (but NOT yaw):
\begin{align}
h(\mathbf{x}) = \begin{bmatrix} \phi \\ \theta \end{bmatrix} = \begin{bmatrix} \arctan2(a_y, a_z) \\ \arctan2(-a_x, \sqrt{a_y^2 + a_z^2}) \end{bmatrix}
\end{align}

The Jacobian is simple: $H = \begin{bmatrix} 1 & 0 & 0 & 0 & 0 & 0 \\ 0 & 1 & 0 & 0 & 0 & 0 \end{bmatrix}$.

\textbf{Key design decisions:}
\begin{itemize}[nosep]
    \item \textbf{Yaw is unobservable} from a 6-axis IMU alone (gravity does not tell you which direction is north). Yaw relies entirely on the gyroscope. To correct yaw drift, add a magnetometer (9-DOF IMU) or external heading reference.
    \item \textbf{$Q$ matrix tuning:} The diagonal entries for angles ($Q_{00}, Q_{11}, Q_{22}$) represent how much the gyro integration can drift per step. Larger = trust gyro less. The bias entries ($Q_{33}, Q_{44}, Q_{55}$) represent bias random walk rate --- typically very small ($10^{-4}$ to $10^{-6}$).
    \item \textbf{$R$ matrix tuning:} Represents accelerometer noise. During vibration or linear acceleration, the accelerometer readings are corrupted. Increase $R$ during high-acceleration maneuvers (adaptive EKF) to avoid trusting bad accelerometer data.
    \item \textbf{Gimbal lock:} When pitch $= \pm 90$\textdegree{}, the $\tan\theta$ and $1/\cos\theta$ terms in the kinematic equations blow up. For robots that can tilt past $\pm 60$\textdegree{}, use \textbf{quaternion representation} instead of Euler angles (same EKF structure, different state parameterization).
\end{itemize}
\end{mdframed}

\begin{figure}[H]
\centering
\includegraphics[width=\textwidth]{ekf_imu_6dof.pdf}
\caption{EKF 6-DOF IMU pose estimation with auto-tuned continuous adaptive $R$. Uses smooth exponential scaling $R = R_\text{base} \cdot (1 + \alpha \, e^{\beta \cdot \Delta a})$ with separate $R_\text{base}$ for roll and pitch, plus pitch-specific x-axis accel deviation detection. Parameters found by offline grid search minimizing weighted RMSE. Top row: roll (RMSE 0.17\textdegree), pitch (0.28\textdegree), yaw (1.54\textdegree). During vibration ($t = 5$--$7$s, shaded), the adaptive $R$ scales up exponentially, rejecting corrupted accelerometer data. Yaw drift is bounded by bias estimation alone (no heading reference). Bottom row: gyroscope biases converge smoothly to true values.}
\end{figure}

\textbf{Beyond EKF:} The EKF's linearization can be inaccurate when the state uncertainty is large or the nonlinearity is strong. Alternatives:
\begin{itemize}
    \item \textbf{Unscented Kalman Filter (UKF):} Instead of linearizing, propagates ``sigma points'' (a small set of carefully chosen sample states) through the nonlinear function. More accurate than EKF for highly nonlinear systems, slightly more expensive (2$n$+1 function evaluations per step where $n$ is the state dimension).
    \item \textbf{Error-State Kalman Filter (ESKF):} Estimates a \textit{small error} around a nominal trajectory instead of the full state. Popular for IMU-based systems because the error state stays small (linearization is accurate) and quaternion singularities are avoided. This is what most production-grade IMU fusion libraries use internally.
    \item \textbf{Complementary filter:} For simple roll/pitch estimation where computational resources are extremely limited, the complementary filter (Section 2) remains a viable one-line alternative.
\end{itemize}

\subsubsection{Quaternion-Based EKF: Avoiding Gimbal Lock}

The Euler-angle EKF above works well when pitch stays within $\pm 60$\textdegree{}. But when pitch approaches $\pm 90$\textdegree{}, the kinematic equations blow up: $\tan\theta \to \infty$ and $1/\cos\theta \to \infty$. This is \textbf{gimbal lock}---a mathematical singularity, not a physical one. The robot is fine; your math is broken.

The fix is to represent orientation with a \textbf{unit quaternion} $\mathbf{q} = \begin{bmatrix} q_w & q_x & q_y & q_z \end{bmatrix}^T$ instead of Euler angles. Quaternions have no singularities, are compact (4 numbers vs 9 for a rotation matrix), and compose rotations efficiently.

\textbf{State vector:} $\mathbf{x} = \begin{bmatrix} \mathbf{q}^T & \boldsymbol{\beta}^T \end{bmatrix}^T \in \mathbb{R}^7$ where $\boldsymbol{\beta} = \begin{bmatrix} \beta_x & \beta_y & \beta_z \end{bmatrix}^T$ are the gyroscope biases.

\textbf{Prediction step:} Given gyroscope measurement $\boldsymbol{\omega}_m$, the bias-corrected angular velocity is $\boldsymbol{\omega} = \boldsymbol{\omega}_m - \boldsymbol{\beta}$. The quaternion kinematics are:
\begin{equation}
\dot{\mathbf{q}} = \frac{1}{2}\mathbf{q} \otimes \begin{bmatrix} 0 \\ \boldsymbol{\omega} \end{bmatrix}
= \frac{1}{2}\Omega(\boldsymbol{\omega})\,\mathbf{q}
\end{equation}

where $\otimes$ is quaternion multiplication and $\Omega(\boldsymbol{\omega})$ is the $4\times 4$ skew matrix:
\begin{equation}
\Omega(\boldsymbol{\omega}) = \begin{bmatrix}
0 & -\omega_x & -\omega_y & -\omega_z \\
\omega_x & 0 & \omega_z & -\omega_y \\
\omega_y & -\omega_z & 0 & \omega_x \\
\omega_z & \omega_y & -\omega_x & 0
\end{bmatrix}
\end{equation}

Discrete-time prediction (first-order):
\begin{equation}
\mathbf{q}_{k+1} = \left(I + \frac{\Delta t}{2}\Omega(\boldsymbol{\omega}_k)\right)\mathbf{q}_k, \qquad
\boldsymbol{\beta}_{k+1} = \boldsymbol{\beta}_k
\end{equation}

After prediction, \textbf{normalize} $\mathbf{q}$ to unit length: $\mathbf{q} \leftarrow \mathbf{q}/\|\mathbf{q}\|$.

\textbf{Measurement model:} The accelerometer measures gravity in the body frame. Given quaternion $\mathbf{q}$, the predicted gravity direction in body frame is:
\begin{equation}
\mathbf{g}_{\text{pred}} = R(\mathbf{q})^T \begin{bmatrix} 0 \\ 0 \\ g \end{bmatrix}
= g\begin{bmatrix}
2(q_x q_z - q_w q_y) \\
2(q_y q_z + q_w q_x) \\
q_w^2 - q_x^2 - q_y^2 + q_z^2
\end{bmatrix}
\end{equation}

The measurement Jacobian $H = \partial \mathbf{g}_{\text{pred}} / \partial \mathbf{x}$ is a $3\times 7$ matrix (3 accel components, 4 quaternion + 3 bias states). The bias columns are zero since accelerometer readings do not depend on gyro biases directly.

\textbf{Update step:} Standard EKF update applies. The innovation is $\boldsymbol{\nu} = \mathbf{a}_{\text{meas}} - \mathbf{g}_{\text{pred}}$, and after the update, normalize $\mathbf{q}$ again.

\begin{mdframed}[linewidth=1pt, innertopmargin=10pt, innerbottommargin=10pt]
\textbf{Euler Angles vs Quaternions: When to Use Which}

\vspace{0.3em}
\begin{itemize}[nosep]
    \item \textbf{Euler angles} are fine if your robot operates within $\pm 60$\textdegree{} pitch (ground vehicles, most turrets, most drones in normal flight). The math is more intuitive and easier to debug.
    \item \textbf{Quaternions} are necessary for robots that can flip, tumble, or point straight up/down (acrobatic drones, underwater ROVs, robot arms near wrist singularities).
    \item \textbf{Tip:} You can use quaternions internally for all computation and convert to Euler angles only for display and human-readable logging: $\phi = \text{atan2}(2(q_w q_x + q_y q_z),\; 1 - 2(q_x^2 + q_y^2))$, etc.
    \item \textbf{The Error-State approach} (ESKF) avoids working with the 4-element quaternion directly in the EKF state. Instead, it estimates a 3-element rotation error $\delta\boldsymbol{\theta}$ and applies it to the quaternion after each update. This keeps the covariance matrix $P$ at $6\times 6$ instead of $7\times 7$ and avoids the unit-norm constraint issue. Most production IMU libraries (e.g., PX4's EKF2, VINS-Mono) use ESKF internally.
\end{itemize}
\end{mdframed}

\subsubsection{Adaptive EKF: Auto-Tuning \texorpdfstring{$Q$}{Q} and \texorpdfstring{$R$}{R} Online}

The EKF's performance depends critically on the noise covariances $Q$ (process noise) and $R$ (measurement noise). In practice, these are often tuned by hand --- a tedious trial-and-error process. Worse, the true noise characteristics \textit{change during operation}: an IMU on a vibrating chassis has very different noise from the same IMU sitting still. Fixed $Q$ and $R$ cannot handle this.

\textbf{Adaptive EKF} methods estimate or adjust $Q$ and $R$ \textit{online}, automatically. There are several approaches, ranging from simple to sophisticated.

\textbf{Method 1: Innovation-Based Adaptive Estimation (IAE)}

This is the simplest and most practical method. The idea: the \textbf{innovation sequence} $\boldsymbol{\nu}_k = \mathbf{y}_k - h(\hat{\mathbf{x}}_{k|k-1})$ (the difference between what you measured and what you predicted) contains information about whether your $Q$ and $R$ are correct.

\textit{The key insight:} If the filter is tuned correctly, the innovation should be \textbf{zero-mean white noise} with covariance $S_k = H_k P_{k|k-1} H_k^T + R$. If the innovations are consistently larger than expected, either $Q$ is too small (you are trusting the model too much) or $R$ is too small (you are trusting the sensors too much).

\textbf{Algorithm:} Estimate the actual innovation covariance from a sliding window of $N_w$ recent innovations:
\begin{equation}
\hat{S}_k = \frac{1}{N_w} \sum_{j=k-N_w+1}^{k} \boldsymbol{\nu}_j \boldsymbol{\nu}_j^T
\end{equation}

Then update $R$ to match:
\begin{equation}
\hat{R}_k = \hat{S}_k - H_k P_{k|k-1} H_k^T
\end{equation}

If $\hat{R}_k$ has negative eigenvalues (which can happen with small $N_w$), clamp them to a minimum value. Typical window: $N_w = 20$--$100$ samples.

\textbf{Code:}
\begin{verbatim}
    // After computing innovation nu = z - h(x_pred):
    // Accumulate outer product in sliding window
    innov_buffer[buf_idx] = nu;
    buf_idx = (buf_idx + 1) % N_WINDOW;

    if (samples_seen >= N_WINDOW) {
        // Estimate innovation covariance
        Mat<NZ,NZ> S_hat = Mat<NZ,NZ>::zeros();
        for (int j = 0; j < N_WINDOW; j++) {
            S_hat = S_hat + innov_buffer[j] * innov_buffer[j].T();
        }
        S_hat = S_hat * (1.0f / N_WINDOW);

        // Adapt R
        R = S_hat - H * P_pred * H.T();
        // Clamp diagonal to minimum (prevent negative/zero variance)
        for (int i = 0; i < NZ; i++)
            if (R(i,i) < R_min) R(i,i) = R_min;
    }
\end{verbatim}

\textbf{Method 2: Residual-Based $Q$ Adaptation}

While IAE adapts $R$, the process noise $Q$ can be adapted using the \textbf{state residuals} (difference between the predicted and updated state):
\begin{equation}
\hat{Q}_k = \frac{1}{N_w} \sum_{j} (K_j \boldsymbol{\nu}_j)(K_j \boldsymbol{\nu}_j)^T + P_{k|k} - F_k P_{k-1|k-1} F_k^T
\end{equation}

This is more complex and less numerically stable than IAE. In practice, many systems only adapt $R$ and leave $Q$ fixed, because $R$ varies more (sensor noise changes with environment) while $Q$ is dominated by model uncertainty (which is more constant).

\textbf{Method 3: Fading Memory (Exponential Forgetting)}

The simplest ``adaptive'' trick: multiply the predicted covariance by a \textbf{fading factor} $\alpha > 1$:
\begin{equation}
P_{k|k-1} = \alpha \cdot (F_k P_{k-1|k-1} F_k^T + Q)
\end{equation}

This artificially inflates the uncertainty at each step, making the filter ``forget'' old data and trust new measurements more. It is equivalent to making $Q$ larger without explicitly changing it.

\begin{itemize}
    \item $\alpha = 1.0$: standard EKF (no forgetting).
    \item $\alpha = 1.01$--$1.05$: mild adaptation --- the filter reacts faster to sudden changes.
    \item $\alpha > 1.1$: very aggressive --- the filter almost ignores the model and follows measurements. Only use during known disturbances.
\end{itemize}

\textbf{When to use:} When you expect occasional sudden model changes (e.g., a robot arm picks up an unknown load --- the inertia changes instantly). The fading memory prevents the filter from ``locking on'' to outdated dynamics.

\textbf{Code (one line added to standard EKF):}
\begin{verbatim}
    P_pred = F * P * F.T() + Q;
    P_pred = P_pred * alpha;   // <-- fading memory, alpha = 1.01
\end{verbatim}

\textbf{Method 4: Multi-Mode Adaptive EKF (for IMU)}

For competition robots, the most practical approach is a simple \textbf{mode-switching} strategy rather than continuous adaptation:

\begin{verbatim}
    // Compute acceleration magnitude
    float accel_mag = sqrt(ax*ax + ay*ay + az*az);
    float accel_err = fabs(accel_mag - 9.81f);

    if (accel_err < 0.5f) {
        // Stationary or slow motion: trust accelerometer
        ekf.R(0,0) = 0.05f;  ekf.R(1,1) = 0.05f;
    } else if (accel_err < 2.0f) {
        // Moderate motion: reduce trust
        ekf.R(0,0) = 0.5f;   ekf.R(1,1) = 0.5f;
    } else {
        // High acceleration / vibration: mostly ignore accelerometer
        ekf.R(0,0) = 10.0f;  ekf.R(1,1) = 10.0f;
    }
\end{verbatim}

\textbf{The logic:} When $|\mathbf{a}| \approx g$ (9.81 m/s$^2$), the robot is not accelerating and the accelerometer truly measures gravity --- trust it ($R$ small). When $|\mathbf{a}|$ deviates from $g$, there is linear acceleration or vibration contaminating the reading --- distrust it ($R$ large).

This is not mathematically elegant, but it is \textbf{extremely effective} in practice. Many championship-winning robots use exactly this approach. It is robust, easy to tune (just three thresholds), and computationally free.

\textbf{Summary: which method to use?}

\begin{center}
\begin{tabular}{llll}
\toprule
\textbf{Method} & \textbf{Adapts} & \textbf{Complexity} & \textbf{Best for} \\
\midrule
Multi-mode switching & $R$ (discrete) & Trivial & IMU on mobile robots \\
Fading memory & Effective $Q$ & 1 line & Sudden model changes \\
IAE (innovation-based) & $R$ (continuous) & Moderate & Varying sensor noise \\
Residual-based & $Q$ and $R$ & High & Research / offline \\
\bottomrule
\end{tabular}
\end{center}

\textbf{A word of caution:} Adaptive methods can \textit{destabilize} the filter if implemented incorrectly. Always clamp adapted $Q$ and $R$ to reasonable ranges, use a minimum window size ($N_w \geq 20$), and validate on recorded data before deploying on hardware.

\subsubsection{Luenberger Observer vs.\ Kalman Filter: A Comprehensive Comparison}

Now that you know both observers, let us put them side by side and understand when to use which.

\begin{center}
\begin{tabular}{p{3.5cm}p{5.5cm}p{5.5cm}}
\toprule
\textbf{Aspect} & \textbf{Luenberger Observer} & \textbf{Kalman Filter} \\
\midrule
\textbf{Design input} & Desired pole locations (you choose) & Noise covariances $Q_n$, $R_n$ (from physics or tuning) \\
\textbf{Gain computation} & One-time pole placement & Recursive (updates $P$ and $K_k$ every step) or steady-state \\
\textbf{Optimality} & Not optimal --- you pick ``reasonable'' poles & Optimal (minimizes mean-squared error for linear Gaussian systems) \\
\textbf{Noise handling} & No explicit noise model --- gain is fixed & Explicitly models process and measurement noise \\
\textbf{Adaptivity} & Fixed gain $L$ --- same in all conditions & Gain $K_k$ adapts as uncertainty ($P$) evolves \\
\textbf{Computational cost} & Very low (matrix multiply only) & Higher (matrix multiply + inverse each step) \\
\textbf{Tuning difficulty} & Medium (choose pole locations) & Medium (choose $Q_n$, $R_n$) \\
\textbf{Confidence output} & None & $P$ matrix tells you how confident the estimate is \\
\textbf{Sensor fusion} & Possible but manual & Natural --- just expand $C$ matrix with more sensors \\
\bottomrule
\end{tabular}
\end{center}

\textbf{The deep connection:} At steady state ($k \to \infty$), the Kalman filter gain $K_k$ converges to a constant matrix $K_\infty$. At that point, the Kalman filter \textit{is} a Luenberger observer with $L = K_\infty$. The difference is that the Kalman filter \textit{computed} this $L$ optimally from the noise statistics, while the Luenberger approach requires you to choose $L$ by intuition or pole placement.

In fact, the steady-state Kalman gain $K_\infty$ is the solution to the observer Riccati equation --- the exact dual of the controller Riccati equation that LQR solves. This completes the beautiful symmetry:

\begin{center}
\begin{tabular}{lcc}
\toprule
 & \textbf{Control (input) side} & \textbf{Estimation (output) side} \\
\midrule
Manual design & Pole placement for $K$ & Luenberger: pole placement for $L$ \\
Optimal design & LQR: minimize $\int (\mathbf{x}^TQ\mathbf{x} + \mathbf{u}^TR\mathbf{u})\,dt$ & Kalman: minimize $E[\|\mathbf{x} - \hat{\mathbf{x}}\|^2]$ \\
Riccati variable & $P$ (control cost-to-go) & $P$ (estimation uncertainty) \\
Tuning knobs & $Q$ vs.\ $R$ (state penalty vs.\ effort) & $Q_n$ vs.\ $R_n$ (model trust vs.\ sensor trust) \\
\bottomrule
\end{tabular}
\end{center}

\subsubsection{LQG (Linear Quadratic Gaussian)}

Now that we have both LQR (optimal controller) and the Kalman filter (optimal observer), we can combine them. \textbf{LQG (Linear Quadratic Gaussian)} is what happens when you do: LQR handles control, Kalman handles estimation, and the separation principle guarantees the combination is stable.

\textbf{The LQG architecture:}
\begin{align}
\text{Kalman filter:} \quad \dot{\hat{\mathbf{x}}} &= A\hat{\mathbf{x}} + B\mathbf{u} + L(\mathbf{y} - C\hat{\mathbf{x}}) \\
\text{LQR controller:} \quad \mathbf{u} &= -K\hat{\mathbf{x}}
\end{align}

where $K$ comes from the control Riccati (minimizing state cost) and $L$ comes from the estimation Riccati (minimizing estimation error). The controller never sees the true state --- it acts on the Kalman estimate.

\textbf{Why LQG is superior to PID for MIMO/noisy systems:}
\begin{enumerate}
    \item \textbf{Optimal noise handling:} PID amplifies noise through the derivative term. LQG's Kalman filter provides the minimum-variance state estimate --- the cleanest possible signal for the controller.
    \item \textbf{MIMO natural:} PID is fundamentally SISO. For a 4-state balance bot, PID needs cascaded loops with hand-tuned interactions. LQG handles all states in one $K$ matrix.
    \item \textbf{State estimation:} PID cannot estimate unmeasured states. LQG's Kalman filter reconstructs the full state from partial measurements (e.g., estimating velocity from position-only sensors).
    \item \textbf{Systematic tuning:} PID has 3 knobs per loop. LQG has $Q$, $R$ (control cost) and $Q_n$, $R_n$ (noise model) --- physically meaningful matrices.
\end{enumerate}

\textbf{The LQG robustness caveat:} Pure LQR has guaranteed stability margins ($\geq$ 6 dB gain, $\geq$ 60\textdegree{} phase). When you add the Kalman filter (making it LQG), these guarantees \textit{vanish}. The LQG controller can have arbitrarily small stability margins. This is the famous \textbf{Doyle's counter-example} (1978) --- it motivated the development of $H_\infty$ robust control.

\textbf{Practical implication:} LQG works well when your model is accurate and noise is well-characterized. If there is significant model uncertainty (unknown friction, varying load), consider:
\begin{itemize}
    \item \textbf{LQG/LTR (Loop Transfer Recovery):} Modify the Kalman filter tuning to recover LQR's robustness margins at the plant input.
    \item \textbf{$H_\infty$ control:} Design for worst-case model uncertainty instead of average-case noise.
    \item \textbf{Practical approach:} Validate LQG in simulation with $\pm 30\%$ parameter variations. If it stays stable, it is usually fine for competition use.
\end{itemize}

\begin{figure}[H]
\centering
\includegraphics[width=\textwidth]{lqg_comparison.pdf}
\caption{LQG (blue) vs direct noisy feedback (red) on a balance bot with disturbance at $t=2$s. Both use the same LQR gains, but red uses raw noisy measurements + numerical differentiation while blue uses a Kalman filter. LQG recovers faster (top-left), maintains position (top-right), uses smoother control (bottom-left), and achieves lower RMS error (bottom-right). This demonstrates the value of the Kalman filter: same gains, same structure, dramatically different performance.}
\end{figure}

\textbf{Practical decision tree:}
\begin{enumerate}
    \item \textbf{Do you need state estimation at all?} If you can measure all states directly (e.g., encoder gives position AND velocity), skip the observer entirely.
    \item \textbf{Is your model accurate and noise is not a big concern?} Use Luenberger. It is simpler to implement and debug.
    \item \textbf{Is sensor noise significant or do you need sensor fusion?} Use Kalman. It handles noise optimally and fuses multiple sensors naturally.
    \item \textbf{Is your system nonlinear?} Use Extended Kalman Filter (EKF). It linearizes at each step, which works well for ``mildly'' nonlinear systems like IMU orientation estimation.
    \item \textbf{Is your system highly nonlinear or does EKF diverge?} Use Unscented Kalman Filter (UKF) or consider particle filters (beyond our scope here).
\end{enumerate}

\vspace{0.5em}
\begin{mdframed}[linewidth=1pt, innertopmargin=10pt, innerbottommargin=10pt]
\textbf{Case Study: IMU Tilt Estimation --- Kalman Filter vs.\ Complementary Filter}

\vspace{0.3em}
You are building a self-balancing robot. You need to know the tilt angle $\theta$ accurately and in real time. You have two sensors:

\begin{itemize}[nosep]
    \item \textbf{Accelerometer:} Measures gravity direction $\to$ gives you $\theta_{\text{accel}} = \arctan(a_x / a_z)$. Accurate \textit{on average} but noisy and corrupted by any linear acceleration (when the robot moves, it ``sees'' phantom tilt).
    \item \textbf{Gyroscope:} Measures angular rate $\dot{\theta}_{\text{gyro}}$. Very clean in the short term, but \textit{drifts} over time because you are integrating: $\theta_{\text{gyro}}(t) = \theta_0 + \int \dot{\theta}_{\text{gyro}} \, dt$. A tiny bias of 0.1\textdegree/s becomes 6\textdegree{} of error after one minute.
\end{itemize}

\textbf{The insight:} The accelerometer is good at low frequencies (static tilt) but bad at high frequencies (vibration, acceleration). The gyroscope is good at high frequencies (fast rotation) but bad at low frequencies (drift). They are \textit{complementary}.

\textbf{Solution 1: Complementary filter (the quick and dirty way)}
\[
\theta[n] = \alpha \cdot (\theta[n-1] + \dot{\theta}_{\text{gyro}} \cdot T_s) + (1 - \alpha) \cdot \theta_{\text{accel}}
\]

This is literally a high-pass filter on the gyroscope and a low-pass filter on the accelerometer, with $\alpha \approx 0.98$. It works surprisingly well for a single line of code.

\textbf{Solution 2: Kalman filter (the proper way)}

State: $\mathbf{x} = \begin{bmatrix} \theta \\ \dot{\theta}_{\text{bias}} \end{bmatrix}$ (angle and gyro bias).

Model: the angle changes by the (bias-corrected) gyro rate:
\[
\begin{bmatrix} \theta \\ \dot{\theta}_{\text{bias}} \end{bmatrix}_{k+1} =
\begin{bmatrix} 1 & -T_s \\ 0 & 1 \end{bmatrix}
\begin{bmatrix} \theta \\ \dot{\theta}_{\text{bias}} \end{bmatrix}_k +
\begin{bmatrix} T_s \\ 0 \end{bmatrix} \dot{\theta}_{\text{gyro},k}
\]

Measurement: the accelerometer gives us a noisy observation of $\theta$:
\[
y_k = \begin{bmatrix} 1 & 0 \end{bmatrix} \mathbf{x}_k + v_k
\]

\textbf{What the Kalman filter does that the complementary filter cannot:}
\begin{itemize}[nosep]
    \item Automatically \textbf{estimates and corrects gyro bias} ($\dot{\theta}_{\text{bias}}$). This eliminates drift without manual calibration.
    \item Adapts the \textbf{trust ratio dynamically} based on noise statistics, rather than using a fixed $\alpha$.
    \item Provides a \textbf{confidence estimate} ($P$ matrix) for the state, so you know how much to trust the output.
    \item Generalizes easily to 3D orientation (roll, pitch, yaw) by expanding the state vector.
\end{itemize}

\textbf{When is the complementary filter ``good enough''?} For a simple balance bot or a gimbal where the robot does not accelerate much. The complementary filter is one line of code and requires no matrix math.

\textbf{When do you need the Kalman filter?} When the robot undergoes significant linear acceleration (sentry robot driving fast, drone), when you need gyro bias estimation, or when fusing more than two sensors (add magnetometer for yaw, add wheel encoder for velocity).

\textbf{Competition tip:} Start with the complementary filter. If you see drift or poor performance during aggressive maneuvers, upgrade to a Kalman filter. The complementary filter is also a great ``sanity check''---if your Kalman filter gives a wildly different answer, something is wrong with your model.
\end{mdframed}

\begin{figure}[H]
\centering
\includegraphics[width=\textwidth]{kalman_filter.pdf}
\caption{Top: IMU tilt estimation comparison. Pure gyro integration (orange) drifts hopelessly. The accelerometer (red, faint) is noisy. Both the complementary filter (purple) and Kalman filter (blue) track the true angle (black dashed), but Kalman is smoother. Bottom: the Kalman filter automatically estimates and converges to the true gyro bias of 1.5\textdegree/s.}
\end{figure}

\vspace{0.5em}
\begin{mdframed}[linewidth=1pt, innertopmargin=10pt, innerbottommargin=10pt]
\textbf{Case Study: LQR Balance Bot --- From Equations to Riding}

\vspace{0.3em}
Let us design the controller for a two-wheeled self-balancing robot (like a Segway). This is the ``hello world'' of modern control, and it beautifully showcases why LQR + Kalman is more natural than PID for MIMO systems.

\textbf{The system:} A two-wheeled platform with a tall body that must balance upright while also tracking a desired position.

\textbf{States:} $\mathbf{x} = \begin{bmatrix} \theta & \dot{\theta} & x & \dot{x} \end{bmatrix}^T$ where $\theta$ is tilt angle, $x$ is wheel position.

\textbf{The MIMO challenge:} To move forward, the robot must first tilt \textit{forward} (like a human leaning to walk). But tilting is the thing we are trying to prevent! Position control and balance are \textbf{coupled and conflicting}. A PID approach would need at least two nested loops with careful cross-coupling management. LQR handles this with a single gain matrix $K = [k_1, k_2, k_3, k_4]$:
\[
u = -k_1 \theta - k_2 \dot{\theta} - k_3 (x - x_{\text{ref}}) - k_4 \dot{x}
\]

\textbf{Tuning via Q and R:}
\begin{itemize}[nosep]
    \item Want aggressive balancing? Increase $Q_{11}$ (penalize tilt angle).
    \item Want tight position tracking? Increase $Q_{33}$ (penalize position error).
    \item Motors are weak? Increase $R$ (penalize control effort).
\end{itemize}

A typical starting point: $Q = \text{diag}(100, 1, 10, 1)$, $R = 1$. This says: ``I care 100$\times$ more about tilt than velocity, and 10$\times$ more about position than velocity.'' Run \texttt{K = lqr(A, B, Q, R)} in MATLAB, upload the four gains, and the robot balances.

\textbf{The complete system:}
\begin{enumerate}[nosep]
    \item \textbf{Kalman filter} estimates $[\theta, \dot{\theta}, x, \dot{x}]$ from IMU + wheel encoders.
    \item \textbf{LQR controller} computes motor torque from the estimated state.
    \item Both run at 500--1000 Hz on an STM32.
\end{enumerate}

This Kalman+LQR architecture is the gold standard for balance bots, drone attitude control, and any system where multiple states are coupled.
\end{mdframed}

\begin{figure}[H]
\centering
\includegraphics[width=\textwidth]{lqr_balance.pdf}
\caption{LQR balance bot simulation. Starting from an 8.6\textdegree{} tilt, the controller recovers all four states to zero. Top-left: tilt angle converges in $\sim$1s. Top-right: angular rate. Bottom-left: wheels move backward to catch the falling body, then return. Bottom-right: control force --- note the initial burst followed by smooth decay.}
\end{figure}

\newpage

\subsection{Trajectory Generation: Telling Your Motor Where to Go}

You cannot just command a step from $0^\circ$ to $90^\circ$ and expect good things to happen. The motor cannot teleport. You need a smooth \textbf{reference trajectory}---a path through position, velocity, and acceleration that respects the physical limits of your actuator. Without one, your PID saturates immediately, integral windup kicks in, and by the time the motor gets anywhere near $90^\circ$ the controller is fighting itself. Trajectory generation is what separates ``my turret kind of works'' from ``my turret tracks like a top team's.''

\subsubsection{Trapezoidal Velocity Profile}

The trapezoidal profile is the workhorse of motion control. Given a move of distance $d$ with maximum velocity $v_{\max}$ and maximum acceleration $a_{\max}$, it produces three phases:

\begin{enumerate}[nosep]
    \item \textbf{Accelerate} at $a_{\max}$ from rest up to $v_{\max}$
    \item \textbf{Cruise} at $v_{\max}$
    \item \textbf{Decelerate} at $-a_{\max}$ back to rest
\end{enumerate}

\textbf{Timing and distances:}
\begin{align}
t_{\text{accel}} &= \frac{v_{\max}}{a_{\max}}, \qquad
d_{\text{accel}} = \frac{1}{2}a_{\max}t_{\text{accel}}^2, \qquad
d_{\text{cruise}} = d - 2\,d_{\text{accel}}
\end{align}

If $d_{\text{cruise}} < 0$, there is not enough distance to reach $v_{\max}$. The profile becomes \textbf{triangular}: the motor accelerates to a peak velocity $v_{\text{peak}} = \sqrt{a_{\max}\,d}$ and immediately decelerates. No cruise phase.

\textbf{Reference signals at time $t$:}

Let $t_1 = t_{\text{accel}}$, $t_2 = t_1 + d_{\text{cruise}}/v_{\max}$, $t_3 = t_2 + t_{\text{accel}}$.

\begin{equation}
\theta_{\text{ref}}(t) = \begin{cases}
\tfrac{1}{2}a_{\max}t^2 & 0 \le t < t_1 \quad \text{(accelerate)}\\[4pt]
d_{\text{accel}} + v_{\max}(t - t_1) & t_1 \le t < t_2 \quad \text{(cruise)}\\[4pt]
d - \tfrac{1}{2}a_{\max}(t_3-t)^2 & t_2 \le t \le t_3 \quad \text{(decelerate)}
\end{cases}
\end{equation}

\begin{equation}
\dot{\theta}_{\text{ref}}(t) = \begin{cases}
a_{\max}\,t & 0 \le t < t_1 \\
v_{\max} & t_1 \le t < t_2 \\
a_{\max}(t_3 - t) & t_2 \le t \le t_3
\end{cases}
\qquad
\ddot{\theta}_{\text{ref}}(t) = \begin{cases}
+a_{\max} & 0 \le t < t_1 \\
0 & t_1 \le t < t_2 \\
-a_{\max} & t_2 \le t \le t_3
\end{cases}
\end{equation}

These three signals are all you need for a fully feedforward-augmented controller (more on that in the practical box below).

\subsubsection{S-Curve Profile}

The trapezoidal profile has discontinuous acceleration---it jumps from $0$ to $a_{\max}$ instantaneously. That sudden jerk excites mechanical resonances, stresses gears, and causes vibration. The \textbf{S-curve} (or jerk-limited) profile adds a seventh constraint: the \emph{rate of change} of acceleration (jerk, $j_{\max}$) is bounded.

This turns the three-phase trapezoidal into a seven-phase profile:
\begin{enumerate}[nosep]
    \item Jerk up (acceleration increases from 0 to $a_{\max}$)
    \item Constant acceleration at $a_{\max}$
    \item Jerk down (acceleration decreases from $a_{\max}$ to 0)
    \item Cruise at $v_{\max}$
    \item Jerk down (acceleration goes to $-a_{\max}$)
    \item Constant deceleration at $-a_{\max}$
    \item Jerk up (back to zero acceleration)
\end{enumerate}

The result is a position profile that is twice differentiable everywhere---smooth enough that gimbals and precision mechanisms barely feel the motion start. The math is more involved, but in practice you will rarely implement this from scratch: motion control libraries (WPILib's \texttt{TrapezoidProfile}, REV's motion magic, etc.) handle it for you.

\begin{mdframed}[linewidth=1pt, innertopmargin=10pt, innerbottommargin=10pt]
\textbf{Practical: Choosing a Profile}

\vspace{0.3em}
\textbf{Trapezoidal} is right for 90\% of competition use cases:
\begin{itemize}[nosep]
    \item Drive wheels, flywheels, turret slew
    \item Simple to implement and debug---three phases, piecewise formulas
    \item Good enough unless you hear gear chatter or see oscillation on fast moves
\end{itemize}

\textbf{S-curve} when smoothness is non-negotiable:
\begin{itemize}[nosep]
    \item Gimbals and two-axis stabilization platforms
    \item Precision positioning mechanisms (camera sliders, encoder-based arms)
    \item Anywhere gear wear or vibration is a concern
\end{itemize}

\textbf{The full feedforward chain for a trajectory-following turret:}

At every control timestep, your trajectory planner outputs three numbers: $\theta_{\text{ref}}$, $\dot{\theta}_{\text{ref}}$, $\ddot{\theta}_{\text{ref}}$. Use all three:
\[
u = \underbrace{K_p\bigl(\theta_{\text{ref}} - \theta\bigr) + K_i\int(\cdot)\,dt + K_d\dot{e}}_{\text{outer position PID}}
  + \underbrace{k_v\,\dot{\theta}_{\text{ref}}}_{\text{velocity FF}}
  + \underbrace{k_a\,\ddot{\theta}_{\text{ref}}}_{\text{accel FF}}
\]

The acceleration feedforward comes from Newton's second law for rotation: $\tau = J\alpha$, so $k_a \approx J / k_{\tau}$ where $k_\tau$ is the motor's torque constant. The position PID only needs to cancel small tracking errors; the feedforward does the heavy lifting.
\end{mdframed}

\begin{figure}[H]
\centering
\includegraphics[width=\textwidth]{trajectory_profiles.pdf}
\caption{Trapezoidal (left) vs S-curve (right) motion profiles for a 90\textdegree{} move. The trapezoidal profile has discontinuous acceleration (sudden jumps), while the S-curve limits jerk for smoother motion. Both produce position, velocity, and acceleration references for feedforward control.}
\end{figure}

\subsubsection{B\'{e}zier Curves and Spline Interpolation}

While trapezoidal and S-curve profiles handle single point-to-point moves well, many applications require smooth trajectories through multiple waypoints or along shaped paths in 2D/3D space. \textbf{B\'{e}zier curves} and \textbf{cubic splines} are two widely used tools for this purpose.

\paragraph{B\'{e}zier Curves.}
A B\'{e}zier curve of degree $n$ is defined by $n+1$ \textbf{control points} $\mathbf{P}_0, \mathbf{P}_1, \ldots, \mathbf{P}_n$. The curve is parameterized by $t \in [0,1]$:
\[
\mathbf{B}(t) = \sum_{i=0}^{n} \binom{n}{i}(1-t)^{n-i}\,t^{i}\,\mathbf{P}_i
\]
The most common variant is the \textbf{cubic B\'{e}zier} ($n=3$), which uses four control points. The curve passes through $\mathbf{P}_0$ (at $t=0$) and $\mathbf{P}_3$ (at $t=1$), while $\mathbf{P}_1$ and $\mathbf{P}_2$ shape the curve without lying on it. A key property is that the tangent at $\mathbf{P}_0$ points toward $\mathbf{P}_1$, and the tangent at $\mathbf{P}_3$ points from $\mathbf{P}_2$---this gives intuitive control over entry and exit directions. B\'{e}zier curves are well suited for \textbf{waypoint-based path planning} where smooth curves with controllable entry/exit directions are needed (e.g., swerve drive paths, robotic arm end-effector trajectories).

\paragraph{Cubic Spline Interpolation.}
To pass through a sequence of waypoints $\{q_0, q_1, \ldots, q_N\}$, one can use \textbf{piecewise cubic polynomials} joined at the waypoints. Each segment takes the form:
\[
q_i(t) = a_i + b_i\,t + c_i\,t^2 + d_i\,t^3
\]
At each junction, \textbf{continuity conditions} require that position, velocity, and acceleration match between adjacent segments, yielding $C^2$ continuity (continuous up to the second derivative). Two common boundary conditions are:
\begin{itemize}
    \item \textbf{Natural spline}: the second derivative is set to zero at both endpoints.
    \item \textbf{Clamped spline}: the velocity (first derivative) is specified at both endpoints.
\end{itemize}
Compared to B\'{e}zier curves, cubic splines are \textbf{guaranteed to pass through all waypoints}---B\'{e}zier curves only interpolate the first and last control points. Compared to trapezoidal profiles, cubic splines provide \textbf{smooth velocity and acceleration everywhere} with no discontinuities.

\paragraph{When to Use What.}
\begin{itemize}
    \item \textbf{Trapezoidal / S-curve}: single point-to-point moves with explicit velocity and acceleration limits.
    \item \textbf{B\'{e}zier curves}: smooth 2D/3D paths with shape control (e.g., swerve drive paths, end-effector trajectories).
    \item \textbf{Cubic splines}: multi-waypoint trajectories where passing through every point is required (e.g., autonomous navigation sequences).
\end{itemize}

\begin{figure}[H]
\centering
\includegraphics[width=\textwidth]{bezier_spline.pdf}
\caption{Left: a cubic B\'{e}zier curve defined by four control points $P_0$--$P_3$. The curve passes through $P_0$ and $P_3$; $P_1$ and $P_2$ shape the curve. Green and red arrows show the tangent directions at the endpoints. Right: cubic spline interpolation passes through all waypoints (red dots), while a B\'{e}zier curve with the same endpoints only interpolates the first and last points.}
\end{figure}
